% Front matter — équipe LogiFlow

\chapter*{Dédicaces}
\addcontentsline{toc}{chapter}{Dédicaces}

Nous dédions ce travail à nos parents, dont le soutien, les sacrifices et la confiance ont rendu possible notre parcours à l'ENSI Tanger. Leur présence, dans les moments de doute comme dans les moments de réussite, a été le fondement de notre progression.

À nos familles et à nos proches, pour leur patience, leurs encouragements et leur bienveillance tout au long de cette formation.

À nos amis, compagnons de route, pour les échanges, l'entraide et la légèreté qu'ils ont su apporter durant les périodes les plus exigeantes du projet.

Enfin, nous dédions ce rapport à toute l'équipe LogiFlow : ce document est le fruit d'un travail collectif, de décisions partagées et d'un engagement commun pour concevoir un système utile, structuré et professionnel.

\vspace{1.2em}
\begin{flushright}
\textit{L'équipe LogiFlow}
\end{flushright}

\clearpage
\chapter*{Remerciements}
\addcontentsline{toc}{chapter}{Remerciements}

Nous tenons à remercier Allah, Le Tout-Puissant, de nous avoir accordé la force et la persévérance nécessaires pour mener à bien ce projet tutoré.

Nous exprimons notre profonde gratitude à notre encadrant, \textbf{Jaber EL BOUHDIDI}, pour son suivi, ses conseils et sa disponibilité tout au long du projet.

Nos remerciements vont à M.~\textbf{CHAKER Jihad}, responsable de la filière Génie informatique, ainsi qu'à l'ensemble des enseignants de l'École des Nouvelles Sciences et Ingénierie de Tanger pour la qualité de la formation reçue.

Enfin, nous remercions les membres du jury d'avoir accepté d'évaluer ce travail et de nous consacrer leur temps et leur expertise.

\clearpage
\chapter*{Résumé}
\addcontentsline{toc}{chapter}{Résumé}

Ce rapport présente le projet tutoré \textbf{LogiFlow}, un \emph{Transport Management System} (TMS) pour le transport routier de marchandises, réalisé en équipe de quatre étudiants ingénieurs à l'ENSI Tanger (filière Génie informatique, année 2025--2026).

L'objectif est de concevoir et de développer une plateforme web unique, filtrée par rôle métier (exploitant, responsable d'exploitation, atelier, commercial, chauffeur, administrateur), couvrant le cycle commercial $\rightarrow$ dossier de transport $\rightarrow$ voyage $\rightarrow$ suivi, carburant et maintenance.

Le backend est un monolithe modulaire \textbf{Spring Boot} / \textbf{Spring Modulith}, structuré en architecture hexagonale par module métier, avec PostgreSQL. Le frontend est une SPA \textbf{Angular} (signals, Tailwind). L'intelligence artificielle (copilote conversationnel SSE, planification assistée, calcul d'itinéraire, maintenance prédictive) est isolée dans un service \textbf{Flask} (LLM cloud compatible OpenAI, OSRM, pgvector), jamais exposé au navigateur : Spring agit comme BFF. L'infrastructure cloud AWS (une EC2, Caddy, Docker Compose) est provisionnée par \textbf{Terraform} et configurée par \textbf{Ansible}, sans SSH exposé (AWS SSM).

Ce travail a permis de mobiliser des compétences en conception métier, architecture logicielle, développement full-stack, intégration d'agents IA et infrastructure as code.

\vspace{1em}
\noindent\textbf{Mots-clés :} Projet tutoré, TMS, Transport, Spring Modulith, Architecture hexagonale, Angular, Keycloak, Flask, Groq, OSRM, Terraform, Ansible, AWS, Intelligence artificielle.

\clearpage
\chapter*{Abstract}
\addcontentsline{toc}{chapter}{Abstract}

This report presents the supervised team project \textbf{LogiFlow}, a road-freight \emph{Transport Management System} (TMS) developed by four engineering students at ENSI Tangier (Computer Engineering, academic year 2025--2026).

The goal is a single role-filtered web application covering the commercial order $\rightarrow$ transport dossier $\rightarrow$ trip $\rightarrow$ tracking, fuel and maintenance lifecycle.

The backend is a modular \textbf{Spring Boot} / \textbf{Spring Modulith} monolith with hexagonal modules and PostgreSQL. The frontend is an \textbf{Angular} SPA. AI features (SSE copilot, assisted planning, routing, predictive maintenance) run in an isolated \textbf{Flask} service with a cloud OpenAI-compatible LLM and \textbf{OSRM}; the browser never talks to Flask. AWS runs a single EC2 Docker stack behind Caddy, managed with \textbf{Terraform} and \textbf{Ansible} via SSM (no public SSH).

\vspace{1em}
\noindent\textbf{Keywords:} Supervised project, TMS, freight, Spring Modulith, hexagonal architecture, Angular, Keycloak, Flask, Groq, OSRM, Terraform, Ansible, AWS, Artificial Intelligence.

\clearpage
\chapter*{Avant-propos}
\addcontentsline{toc}{chapter}{Avant-propos}

\begin{description}[style=nextline,leftmargin=3.2cm,font=\normalfont\bfseries]
  \item[Établissement] École des Nouvelles Sciences \& Ingénierie (ENSI), Tanger
  \item[Filière] Génie informatique --- temps aménagé
  \item[Cadre] Projet tutoré (travail d'équipe)
  \item[Intitulé] LogiFlow --- Transport Management System
  \item[Encadrant] Jaber EL BOUHDIDI
\end{description}

\vspace{1.2em}
\noindent\textbf{Composition de l'équipe et responsabilités} :

\begin{table}[ht]
\centering
\small
\begin{tabular}{@{}clp{3.2cm}p{5.8cm}@{}}
\toprule
\# & Nom et prénom & Rôle(s) & Responsabilités principales \\
\midrule
1 & STITOU Ibrahim
  & Chef de projet ;\newline Responsable Backend
  & Coordination de l'équipe, suivi du planning, communication avec l'encadrant ;\newline développement API Spring Boot, gestion de la base de données \\[0.7em]
2 & ZOUAINE Oussama
  & Responsable Frontend ;\newline Responsable Infra / DevOps
  & Développement Angular, UI/UX, intégration API ;\newline Docker, CI/CD, déploiement AWS \\[0.7em]
3 & RAHALI Hamza
  & Responsable Sécurité
  & Authentification, gestion des accès, sécurité applicative \\[0.7em]
4 & ES-SEMLLALI Oussama
  & Responsable IA
  & Développement et intégration des fonctionnalités IA dans l'application \\
\bottomrule
\end{tabular}
\end{table}

\clearpage
\chapter*{Liste des sigles et acronymes}
\addcontentsline{toc}{chapter}{Liste des sigles et acronymes}

\begin{table}[ht]
\centering
\caption{Sigles et acronymes utilisés dans ce rapport.}
\label{tab:sigles}
\begin{tabular}{@{}lp{10cm}@{}}
\toprule
Acronyme & Désignation \\
\midrule
ADR & Architecture Decision Record \\
API & Application Programming Interface \\
BFF & Backend for Frontend \\
CI/CD & Continuous Integration / Continuous Delivery \\
DTO & Data Transfer Object \\
GHCR & GitHub Container Registry \\
IAM & Identity and Access Management \\
JWT & JSON Web Token \\
LLM & Large Language Model \\
MFA & Multi-Factor Authentication \\
OIDC & OpenID Connect \\
OSRM & Open Source Routing Machine \\
PKCE & Proof Key for Code Exchange \\
RBAC & Role-Based Access Control \\
SSE & Server-Sent Events \\
SSM & AWS Systems Manager Session Manager \\
TMS & Transport Management System \\
VPC & Virtual Private Cloud \\
\bottomrule
\end{tabular}
\end{table}
